\section{Binary, categorical, and observable interfaces}\label{sec:observable53}
\begin{corollary}[Component oscillation for binary means]\label{thm:classification53}
For a finite continuous binary-mean interface $f_{s,j}:H\to[-1,1]$, with the return hypothesis of Theorem~\ref{thm:radius54}, the exact bounded-width threshold in total variation is
\begin{equation}\label{eq:critical53}
 \epsc=\frac14\max_{s,j,c}\left(\max_{h\in c}f_{s,j}(h)-\min_{h\in c}f_{s,j}(h)\right).
\end{equation}
It is attained by a stationary permutation realization. Below it every fixed width has the occupation bound \eqref{eq:occupation54}.
\end{corollary}
\begin{proof}
In the norm $|\cdot|/2$, the radius of an interval is one quarter of its length. Its midpoint lies in $[-1,1]$. Apply Theorem~\ref{thm:radius54}.
\end{proof}
This recovers the scalar component theorem, including all of its nonuniform quantifiers. The vector theorem does not follow by applying this corollary separately to coordinates, because the coordinate centers need not belong to the output simplex.

\begin{proposition}[A three-outcome threshold]\label{prop:categorical54}
On the circle let $\omega=e^{2\pi i/3}$ and, for $j=0,1,2$, set
\[
 p_j(\phi)=\frac19\left|1+e^{i\phi}\omega^{-j}+e^{2i\phi}\omega^{-2j}\right|^2.
\]
For any dense circle command alphabet satisfying the return condition, the categorical law $p=(p_0,p_1,p_2)$ has critical total-variation error $2/3$. Each separately requested binary event $\{j\}$ has threshold $1/2$.
\end{proposition}
\begin{proof}
Orthogonality of the three characters gives $\sum_jp_j=1$; each entry is nonnegative. At $\phi=2\pi j/3$ the vector is the simplex vertex $e_j$. Hence every center $q\in\Delta_3$ has maximum distance at least
\[
 \max_j\TV(e_j,q)=1-\min_jq_j\ge2/3.
\]
The uniform center has distance at most $2/3$ from every point of $\Delta_3$, by convexity and the values at its vertices. Thus the constrained radius is exactly $2/3$. On the other hand each $p_j$ has range $[0,1]$, so its binary mean $2p_j-1$ has oscillation two. Apply Theorem~\ref{thm:radius54} and Corollary~\ref{thm:classification53}. In particular the categorical radius exceeds half the diameter, which here equals $1/2$.
\end{proof}

\begin{corollary}[Joint measure-once unitary laws]\label{cor:quantum54}
Let unitary commands have compact closure $H$ and satisfy the return condition. For a fixed unit seed $\psi_s$ and a finite measurement $E_{j,l}\ge0$, $\sum_lE_{j,l}=I$, prescribe the law
\[
 F_{s,j}(h)_l=\ip{h\psi_s}{E_{j,l}h\psi_s}.
\]
The critical error is its componentwise simplex radius in total variation. For closure $U(d)$ and one orthonormal-basis measurement, this radius is $1-1/d$.
\end{corollary}
\begin{proof}
The displayed map is a continuous simplex-valued polynomial in realified matrix entries. A closed unitary matrix group has finitely many components, so Theorem~\ref{thm:radius54} applies. For $U(d)$, every unit vector is reachable from the seed, and its squared coordinate magnitudes range over the whole simplex. The simplex vertices force radius at least $1-1/d$, attained by the uniform distribution.
\end{proof}
A query selects one complete measurement, not simultaneous answers to incompatible measurements. Several compatible finite outputs, with specified correlations, instead form one joint categorical law. Language recognition at a cutpoint is a different requirement.

For linear binary interfaces on a compact matrix group, write
\[
 f_{s,j}(h)=q_j^Thx_s\in[-1,1],\quad
 V_{\rm r}=\spanop\{hx_s:h\in H,s\in\mathcal S\},
\]
\[
 V_{\rm n}=\{x\in V_{\rm r}:q_j^Thx=0\text{ for all }h\in H,j\in\mathcal J\}.
\]
Both spaces are invariant. The observable space $V_{\rm o}=V_{\rm r}\cap V_{\rm n}^\perp$ is the orthogonal complement of $V_{\rm n}$ \emph{inside} $V_{\rm r}$, canonically realizing the observable quotient.

\begin{theorem}[Intrinsic zero-error criterion]\label{thm:observable53}
For these linear means, bounded exact clocked width is equivalent to finiteness of the generated group's image in $O(V_{\rm o})$, and to trivial action of $K$ on $V_{\rm o}$. An exact stationary permutation realization then exists.
\end{theorem}
\begin{proof}
Corollary~\ref{thm:classification53} identifies bounded exact width with constancy of all interface means on every component. Under that condition, for $k\in K$ and $g,h\in H$, normality gives
\[
 q_j^Th(k-I)gx_s=q_j^T(hkh^{-1})hgx_s-q_j^Thgx_s=0.
\]
Thus $(k-I)V_{\rm r}\subset V_{\rm n}$. On the invariant orthogonal complement $V_{\rm o}$ the difference also lies in $V_{\rm o}$, hence vanishes. The observable action factors through the finite component group. Conversely, if $K$ acts trivially on $V_{\rm o}$, the same difference is invisible and all means are componentwise constant. A finite image of the dense generated group has the same finite image closure from $H$, whose connected subgroup $K$ acts trivially. The zero-dimensional quotient is included.
\end{proof}

\begin{corollary}[The planar threshold]\label{cor:circle53}
Let an orthogonal planar alphabet contain $I$ and an irrational rotation, with any additional orthogonal letters. For signed coordinate seeds $\pm\rho e_1,\pm\rho e_2$, coordinate queries, and $0<\rho\le1$, the threshold is $\epsc=\rho/2$. At and above it one command label with a fair binary decoder suffices. Below it every fixed width has a finite occupation budget.
\end{corollary}
\begin{proof}
The identity component is $SO(2)$, and every specified mean has range $[-\rho,\rho]$ on each component of $H=SO(2)$ or $O(2)$. All midpoint decoders equal zero, so the seed/component register collapses to one label.
\end{proof}
For the noncommuting rational alphabet
\[
 I,\quad R=\begin{pmatrix}3/5&-4/5\\4/5&3/5\end{pmatrix},\quad
 F=\begin{pmatrix}1&0\\0&-1\end{pmatrix},
\]
one has $FRF=R^{-1}\ne R$. The rotation has infinite order: otherwise $(3+4i)/5$ and its inverse would be algebraic integers, whereas their rational sum $6/5$ is not an integer. This example has the sharp threshold above. We do not infer a badly approximable angle from rational matrix entries.

\begin{proposition}[Independent toral coordinates]\label{prop:torus53}
Suppose a component is identified with the full product torus $(\R/2\pi\Z)^m$, surjectively in independent coordinates, and a mean on it equals
\[
 a_0+\sum_{l=1}^m(a_l\cos\theta_l+b_l\sin\theta_l).
\]
Its contribution to \eqref{eq:critical53} is $\tfrac12\sum_l\sqrt{a_l^2+b_l^2}$.
\end{proposition}
\begin{proof}
Each summand has the symmetric range with amplitude $\sqrt{a_l^2+b_l^2}$, and all extrema are simultaneously attainable in the full product. Divide the resulting oscillation by four.
\end{proof}
On a proper subtorus, extrema must instead be taken subject to its relations. Ambient coordinates do not certify independence. Similarly, invisible directions do not create an obstruction: for $H=\{\diag(R_\theta,1)\}$, seed and query $e_3$ give constant mean one, whereas seed $\rho e_1$ and query $e_1$ have threshold $\rho/2$.

\begin{proposition}[Compact similarity]\label{prop:similarity53}
The classification applies to a finite alphabet in $GL(D,\R)$ whose generated group, including inverses, is uniformly bounded, with the executable return hypothesis retained. Conjugating to an invariant Euclidean metric changes neither output probabilities nor stochastic widths.
\end{proposition}
\begin{proof}
Uniform bounds on matrices and inverses make the closure a compact group. The matrix $Q=\int_Hh^Th\,d\mu(h)$ is positive definite and satisfies $A^TQA=Q$. Hence $Q^{1/2}AQ^{-1/2}$ is orthogonal. Transport the continuous output maps through this conjugacy; for a linear mean transport $x_s$ to $Q^{1/2}x_s$ and $q_j$ to $Q^{-1/2}q_j$. Exact identity products and their word lengths are preserved.
\end{proof}

\subsection{A curved categorical image with distinct component centers}
\begin{proposition}\label{prop:curved55}
Let $H=(\R/2\pi\Z)\times\Z/2\Z$, with an identity letter, an irrational rotation in the first factor, and a toggle of the second factor. Put
\[
 q_0=(\tfrac12,\tfrac13,\tfrac16),\qquad
 q_1=(\tfrac16,\tfrac13,\tfrac12),\qquad t=\tfrac1{12},
\]
and prescribe the three-outcome law
\[
 F(\theta,c)_l=q_{c,l}+t\cos(\theta-2\pi l/3),\qquad l=0,1,2.
\]
Each component image is a proper closed curve in the interior of the simplex, with unique constrained total-variation center $q_c$ and radius $t$. The union of both component images has radius $1/4$. Thus the sharp clocked threshold is $1/12$, and the stationary minimum at that threshold is exactly two labels.
\end{proposition}
\begin{proof}
The three cosines sum to zero and every coordinate is at least $1/12$, so $F$ is a legal law in the simplex interior. The first two independent trigonometric coordinates identify a nondegenerate ellipse in its affine plane. For a zero-sum vector $z\in\R^3$, one has $\frac12\norm{z}_1=\max_l|z_l|$. Each coordinate of the curve ranges over $[q_{c,l}-t,q_{c,l}+t]$. Hence for every center $y\in\Delta_3$,
\[
 \sup_\theta\TV(F(\theta,c),y)
       =t+\max_l|q_{c,l}-y_l|.
\]
This proves uniqueness of $q_c$ and radius $t$. For the union, the least possible value of $\max_{c,l}|q_{c,l}-y_l|$ is $1/6$, forced by the two endpoints in the first coordinate and attained by $y=(1/3,1/3,1/3)$. Its radius is therefore $t+1/6=1/4$. A one-label stationary machine has a constant decoder and cannot reach error $t$; two component labels attain it. The clocked threshold follows from Theorem~\ref{thm:radius54}.
\end{proof}

\begin{remark}
The total-variation diameter of a probability simplex is one, so half its diameter is $1/2$; the three-vertex example above has radius $2/3$. In the unitary corollary, $U(d)$ is connected. These facts explain respectively why coordinate midpoints do not solve the categorical problem and why the full-unitary component is the whole group.
\end{remark}
